\documentclass[a4paper]{arrowhead}

\usepackage[yyyymmdd]{datetime}
\usepackage{etoolbox}
\usepackage[utf8]{inputenc}
\usepackage[T1]{fontenc}
\usepackage{multirow}
\usepackage{longtable}
\usepackage{tikz}
\usetikzlibrary{positioning,shapes.multipart,arrows.meta}

\renewcommand{\dateseparator}{-}
\setlength{\emergencystretch}{3em}
\lstset{language=HTTP,breakatwhitespace=false,columns=fullflexible}

\newcommand{\fparam}[1]{\textit{\textcolor{ArrowheadBlue}{#1}}}

%% Special references
\newcommand{\mref}[2][]{{\textcolor{ArrowheadPurple}{\hyperref[sec:model:#2]{\ifstrempty{#1}{#2}{#1}}}}}
\newcommand{\oref}[1]{{\textcolor{ArrowheadBlue}{\hyperref[sec:operations:#1]{#1}}}}
\newcommand{\pdef}[1]{{\textcolor{ArrowheadGrey}{#1 \label{sec:model:primitives:#1} \label{sec:model:primitives:#1s}}}}
\newcommand{\pref}[1]{{\textcolor{ArrowheadGrey}{\hyperref[sec:model:primitives:#1]{#1}}}}

% as template_idd.tex; the input and output arguments are passed already formatted (\mref or \pref)
\newrobustcmd\osubsection[5]{
  \addtocounter{subsection}{1}
  \addcontentsline{toc}{subsection}{\protect\numberline{\thesubsection}operation \textcolor{ArrowheadBlue}{#1}}
  \subsection*{
    \thesubsection\quad
    #2 \textcolor{ArrowheadPurple}{#3} \\
    \small
    \hspace*{0.07\textwidth}\begin{minipage}{0.1\textwidth}
      \vspace*{1mm}
      Operation: \\
      \notblank{#4}{Input: \\}{}
      \notblank{#5}{Output: \\}{}
    \end{minipage}
    \begin{minipage}{0.825\textwidth}
      \vspace*{1mm}
      \textcolor{ArrowheadBlue}{#1} \\
      \notblank{#4}{#4 \\}{}
      \notblank{#5}{#5 \\}{}
    \end{minipage}
  }
  \label{sec:operations:#1}
}
\newrobustcmd\msubsection[2]{
  \addtocounter{subsection}{1}
  \addcontentsline{toc}{subsection}{\protect\numberline{\thesubsection}#1 \textcolor{ArrowheadPurple}{#2}}
  \subsection*{\thesubsection\quad#1 \textcolor{ArrowheadPurple}{#2}}
  \label{sec:model:#2} \label{sec:model:#2s}
}
%%

\begin{document}
\sloppy

%% Arrowhead Document Properties
\ArrowheadTitle{sysMonitor AMQP/TLS/JSON}
\ArrowheadServiceID{sysMonitor}
\ArrowheadType{Interface Design Description}
\ArrowheadTypeShort{IDD}
\ArrowheadVersion{5.2.2}
\ArrowheadDate{2026-11-09}
\ArrowheadAuthor{Jerker Delsing}
\ArrowheadStatus{RELEASE CANDIDATE}
\ArrowheadContact{jerker.delsing@ltu.se}
\ArrowheadFooter{\href{www.arrowhead.eu}{www.arrowhead.eu}}
\ArrowheadSetup
%%

%% Front Page
\begin{center}
  \vspace*{1cm}
  \huge{\arrowtitle}

  \vspace*{0.2cm}
  \LARGE{\arrowtype}
  \vspace*{1cm}

  \Large{Service ID: \textit{"\arrowid"}}

  \vspace*{\fill}

  % Front Page Image

  \vspace*{1cm}
  \vspace*{\fill}

  % Front Page Abstract
  \begin{abstract}
    This document describes the AMQP 0-9-1 protocol with TLS payload security and JSON payload encoding variant of the \texttt{sysMonitor} service of the LocalCloudConfigurator system (SysMonitor-AMQP-IDD-5.2.2).
  \end{abstract}

  \vspace*{1cm}
\end{center}

\newpage
%%

%% Table of Contents
\tableofcontents
\newpage
%%

\section{Overview}
\label{sec:overview}

This document describes the AMQP binding of the native \texttt{sysMonitor} service of LocalCloudConfigurator.

It's implemented using protocol, encoding as stated in below table:
\begin{table}[ht!]
  \centering
  \begin{tabular}{|l|l|l|}
    \hline
    \rowcolor{gray!33} Profile type & Type & Version \\ \hline
    Transfer protocol & AMQP & 0-9-1 \cite{amqp091} \\ \hline
    Data encryption & TLS & 1.3 \cite{rfc8446}, 1.2 \cite{rfc5246} \\ \hline
    Encoding & JSON & RFC 8259 \cite{rfc8259} \\ \hline
    Compression & N/A & - \\ \hline
    Semantics & N/A & - \\ \hline
    Ontology & N/A & - \\ \hline
  \end{tabular}
  \caption{Communication and semantics details used for the \texttt{sysMonitor} service interface}
  \label{tab:comunication_semantics_profile}
\end{table}

This document provides the Interface Design Description IDD to the \textit{sysMonitor -- Service Description} document \cite{lcc-sysMonitor-sd}.

\subsection{Interface registration}

With \texttt{sysmonitor.enabled=true}, the common sysMonitor publisher of LocalCloudConfigurator registers the service at start-up with interface template \texttt{generic\_\allowbreak{}amqps} (protocol \texttt{amqps}, policy \texttt{CERT\_\allowbreak{}AUTH}; \texttt{generic\_\allowbreak{}amqp} without TLS) and the following properties and metadata:

\begin{lstlisting}[caption={Registered service instance (TestCloud).}]
{
  "serviceDefinitionName": "sysMonitor",
  "version": "1.0.0",
  "metadata": {
    "contract": "SysMonitor-AMQP-IDD-5.2.2",
    "schemaVersion": "1.0",
    "publisherMode": "native",
    "observationTypes": [
      "echo",
      "power",
      "status"
    ],
    "delivery": "persistent, publisher-confirmed"
  },
  "interfaces": [
    {
      "templateName": "generic_amqps",
      "protocol": "amqps",
      "policy": "CERT_AUTH",
      "properties": {
        "accessAddresses": [
          "amqp-broker"
        ],
        "accessPort": 5671,
        "exchange": "arrowhead",
        "baseRoutingKey": "TestCloud.sysmonitor.LocalCloudConfigurator.LocalCloudConfigurator",
        "operations": [
          "echo",
          "power",
          "status"
        ]
      }
    }
  ]
}
\end{lstlisting}

\section{Service Operations}
\label{sec:operations}

The publisher is the \texttt{SysMonitor\allowbreak{}Publisher} of the vendored Eclipse Arrowhead common library (RabbitMQ Java client), configured with the \texttt{sysmonitor.\allowbreak{}*} Spring properties (\texttt{enabled}, \texttt{broker.\allowbreak{}address}, \texttt{broker.\allowbreak{}port}, \texttt{local.\allowbreak{}cloud}, \texttt{service.\allowbreak{}instance}, \texttt{publish.\allowbreak{}interval}, \texttt{ssl.\allowbreak{}*}). The echo reports the default health of the common publisher; \texttt{serviceCount} is the number of services in the system information (2: \texttt{general\allowbreak{}Management} and \texttt{local\allowbreak{}Cloud\allowbreak{}Configuration}). The publisher connects to the broker (default \texttt{amqp-\allowbreak{}broker}, port 5671) with TLS, verifies the broker host name and authenticates with its X.509 system certificate (SASL \texttt{EXTERNAL}; the certificate CN is the broker user name). Every publication interval it publishes each observation to the topic exchange \texttt{arrowhead} with the routing key \texttt{<baseRoutingKey>.<operation>}, where the base routing key is \texttt{<Cloud>.sysmonitor.<System>.<serviceInstance>}. Messages are persistent (delivery mode 2), \texttt{application/json}, carry the observation identifier as message id, and are publisher-confirmed (timeout 5 s); a failed cycle is retried with a new connection at the next interval. The broker's topic permissions allow the system to publish only below its own base routing key.

\osubsection{echo}{PUBLISH}{\texttt{\textless{}Cloud\textgreater{}.\allowbreak{}sysmonitor.\allowbreak{}Local\allowbreak{}Cloud\allowbreak{}Configurator.\allowbreak{}Local\allowbreak{}Cloud\allowbreak{}Configurator.\allowbreak{}echo}}{}{\mref{Observation}}

Liveness: \texttt{status} (\texttt{alive} or \texttt{degraded}), \texttt{uptimeSeconds} and a \texttt{detail} text; quality \texttt{good} or \texttt{uncertain}.

\osubsection{power}{PUBLISH}{\texttt{\textless{}Cloud\textgreater{}.\allowbreak{}sysmonitor.\allowbreak{}Local\allowbreak{}Cloud\allowbreak{}Configurator.\allowbreak{}Local\allowbreak{}Cloud\allowbreak{}Configurator.\allowbreak{}power}}{}{\mref{Observation}}

Power supply: \texttt{supplyType=\allowbreak{}line}, \texttt{state=\allowbreak{}available}; unit \texttt{1}.

\osubsection{status}{PUBLISH}{\texttt{\textless{}Cloud\textgreater{}.\allowbreak{}sysmonitor.\allowbreak{}Local\allowbreak{}Cloud\allowbreak{}Configurator.\allowbreak{}Local\allowbreak{}Cloud\allowbreak{}Configurator.\allowbreak{}status}}{}{\mref{Observation}}

Process state: \texttt{state=\allowbreak{}running}, \texttt{health} (\texttt{healthy}/\texttt{degraded}) and \texttt{serviceCount}, the number of registered services.

\begin{lstlisting}[caption={An \oref{echo} observation.}]
{
  "schemaVersion": "1.0",
  "observationId": "LocalCloudConfigurator-echo-1792670400000-42",
  "providerSystem": "LocalCloudConfigurator",
  "serviceInstance": "LocalCloudConfigurator",
  "observationType": "echo",
  "sourceTime": "2026-10-22T12:00:00.000Z",
  "sequence": 42,
  "value": {
    "status": "alive",
    "uptimeSeconds": 3600,
    "detail": "LocalCloudConfigurator is running"
  },
  "quality": "good",
  "metadata": {
    "publisherMode": "native"
  }
}
\end{lstlisting}

\newpage

\section{Data Models}
\label{sec:model}

\msubsection{struct}{Observation}

\noindent\begin{tabularx}{\textwidth}{| p{4.25cm} | p{3.5cm} | X |} \hline
\rowcolor{gray!33} Object Field & Value Type & Description \\ \hline
"schemaVersion" & \pref{String} & \texttt{1.0}. \\ \hline
"observationId" & \pref{String} & \texttt{<System>-<type>-<epoch ms>-<sequence>}. \\ \hline
"providerSystem" & \pref{String} & The publishing system. \\ \hline
"serviceInstance" & \pref{String} & Service instance, default the system name. \\ \hline
"observationType" & \pref{String} & \texttt{echo}, \texttt{power}, \texttt{status}. \\ \hline
"sourceTime" & \pref{String} & ISO 8601 UTC time with milliseconds. \\ \hline
"sequence" & \pref{Number} & Sequence number per publisher. \\ \hline
"value" & \pref{Object} & Observed value, see the operations. \\ \hline
"unit" & \pref{String} & \texttt{1} for \texttt{power} and \texttt{status}; absent for \texttt{echo}. \\ \hline
"quality" & \pref{String} & \texttt{good} or \texttt{uncertain}. \\ \hline
"metadata" & \pref{Object} & \texttt{\{"publisherMode": "native"\}}. \\ \hline
\end{tabularx}

\subsection{Primitives}
\label{sec:model:primitives}

\noindent\begin{tabularx}{\textwidth}{| p{3cm} | X |} \hline
\rowcolor{gray!33} Type & Description \\ \hline
\pdef{Number} & JSON Number, an integer. \\ \hline
\pdef{Object} & JSON Object; \pref{Object}$<$A$>$ has string keys and values of type A. \\ \hline
\pdef{String} & JSON String. \\ \hline
\end{tabularx}

\newpage

\bibliographystyle{IEEEtran}
\bibliography{bibliography}

\newpage

\section{Revision History}
\subsection{Amendments}

\noindent\begin{tabularx}{\textwidth}{| p{1cm} | p{3cm} | p{2cm} | X | p{4cm} |} \hline
\rowcolor{gray!33} No. & Date & Version & Subject of Amendments & Author \\ \hline

1 & 2026-11-09 & \arrowversion & Document created for the sysMonitor service of the Java-Spring LocalCloudConfigurator (common-library publisher). & Claude Opus 5.5 (draft) \\ \hline
\end{tabularx}

\subsection{Quality Assurance}

\noindent\begin{tabularx}{\textwidth}{| p{1cm} | p{3cm} | p{2cm} | X |} \hline
\rowcolor{gray!33} No. & Date & Version & Approved by \\ \hline

1 & & \arrowversion  &  \\ \hline

\end{tabularx}

\end{document}
